\scanpage{162}
\margin{\diagram{\begin{tikzcd}[ampersand replacement=\&,column sep=small,row sep=small]
\&\unsure{\overline{T}|T}\times_Z T'\arrow[dl]\arrow[dr]\&\\
T\arrow[dr]\&\&T'\arrow[dl]\\
\&X\times_S Y=Z\&
\end{tikzcd}}\quad\unsure{f}\quad *}
À vrai dire, \deleted{\unsure{et}} si $\overline{T}$ est l'image de $T$ dans $Z=X\times_S Y$, on voit que $\overline{\alpha^0}$ ``est'' dans \deleted{l'image}
\[
\mathbb{R}^0\Gamma_{\overline{T}}(\mathbb{R}\mathcal{H}om(\mathrm{pr}_1^*(F),\mathrm{pr}_2^!(G_\bullet))
=\mathbb{R}^0\Gamma_{X\times_S Y}(\mathbb{R}^0\underline{\Gamma}_{\overline{T}}\mathbb{R}\mathcal{H}om(\,,))
\]
et de même $\overline{\beta^0}$ ``est \deleted{\unsure{dans l'image}}\inserted{même}'' dans
\[
\mathbb{R}^0\Gamma_{\overline{T}'}(\mathbb{R}\mathcal{H}om(\mathrm{pr}_2^*(G),\mathrm{pr}_1^!(F))
=\mathbb{R}^0\Gamma_{X\times_S Y}(\mathbb{R}^0\underline{\Gamma}_{\overline{T}'},\mathbb{R}\mathcal{H}om
\]
Ainsi, on peut définir un cup-produit
\[
\overline{\alpha^0}\cup\overline{\beta^0}\in
\mathbb{R}^0\Gamma_{\overline{T}\cap\overline{T}'}K_{\unsure{\infty}X\times_S Y}^{\bullet}
\simeq\mathbb{R}^0\Gamma_{\overline{T}\cap\overline{T}'}K_{\overline{T}\times_S\overline{T}'}^{\bullet}
\]
Ceci permet donc de définir, si
\[
\overline{T}\cap\overline{T}'=\bigcup_\alpha S_\alpha,
\]
\deleted{\unsure{réunion}} les $S_\alpha$ ouverts deux à deux disjoints (en nb fini) de définir des composantes canoniques
\[
(\overline{\alpha^0}\cup\overline{\beta^0})_\alpha:
\deleted{\unsure{\in H^0(S_\alpha,K_{S_\alpha}^{\bullet})}}
\]
et des $\deleted{\unsure{\operatorname{Tr}_{S_\alpha}}}\langle\xi^0,\eta^0\rangle_{S_\alpha}$ telles que
\[
\langle\xi^0,\eta^0\rangle=\sum_\alpha\langle\xi^0,\eta^0\rangle_{S_\alpha}
\]
\uline{Exemple} Supposons $\overline{T}\cap\overline{T}'=\varnothing$. Alors le nb de Lefschetz est nul.
\margin{161}
